\documentclass[11pt,a4paper]{article}

% ---------- Datos del documento ----------
\newcommand{\titulo}{Matriz de competencias de la situación de aprendizaje <<Medir lo que no se puede alcanzar>>}
% Título corto para el encabezado de las páginas
\newcommand{\tituloCorto}{Actividad 1}
\newcommand{\asignatura}{Diseño Curricular de las Matemáticas}
\newcommand{\autor}{Javier Gómez López}
% Texto explicativo de la tarea (déjalo vacío para que no aparezca)
\newcommand{\enunciado}{%
Diseñar una situación de aprendizaje para 1.º de Bachillerato con una duración de tres semanas. Como primer paso, elaborar una matriz de competencias que relacione los contenidos que se van a trabajar con los criterios de evaluación de la normativa, operativice cada criterio adaptándolo al contenido y señale los instrumentos de evaluación con los que se evaluará cada criterio operativizado.%
}
\newcommand{\master}{MÁSTER UNIVERSITARIO EN FORMACIÓN DEL PROFESORADO-ESPECIALIDAD EN MATEMÁTICAS -- Universidad Rey Juan Carlos}

% ---------- Paquetes ----------
\usepackage[utf8]{inputenc}
\usepackage[T1]{fontenc}
\usepackage[spanish,es-noshorthands]{babel}
\usepackage{lmodern}
\usepackage[margin=2.5cm,headheight=14pt]{geometry}
\usepackage{xcolor}
\usepackage{titlesec}
\usepackage{fancyhdr}
\usepackage{graphicx}
\usepackage{enumitem}
\usepackage{parskip}
\usepackage{etoolbox}
\usepackage{tabularx}
\usepackage{booktabs}
\usepackage{pdflscape}
\usepackage{tabularray}
\DefTblrTemplate{contfoot-text}{default}{Continúa en la página siguiente}
\DefTblrTemplate{conthead-text}{default}{(continuación)}
\usepackage[most]{tcolorbox}
\usepackage[hidelinks]{hyperref}

% ---------- Estilo ----------
\definecolor{acento}{RGB}{176,0,32}   % rojo URJC
\definecolor{gris}{RGB}{90,90,90}

\titleformat{\section}{\Large\bfseries\color{acento}}{\thesection.}{0.6em}{}
\titleformat{\subsection}{\large\bfseries\color{gris}}{\thesubsection.}{0.6em}{}

\pagestyle{fancy}
\fancyhf{}
\fancyhead[L]{\small\color{gris}\asignatura}
\fancyhead[R]{\small\color{gris}\tituloCorto}
\fancyfoot[C]{\small\thepage}
\renewcommand{\headrule}{{\color{acento}\hrule height 0.6pt}}

\begin{document}

% ---------- Cabecera ----------
\thispagestyle{empty}
\begin{center}
  {\small\color{gris}\MakeUppercase{\master}}\\[1.2em]
  {\color{acento}\rule{\linewidth}{1.2pt}}\\[0.8em]
  {\Huge\bfseries \titulo}\\[0.6em]
  {\Large\color{gris} \asignatura}\\[0.6em]
  {\color{acento}\rule{\linewidth}{1.2pt}}\\[1em]
  {\large \autor}\\[0.3em]
  {\small\color{gris}\today}
\end{center}
\vspace{1.5em}

% ---------- Enunciado ----------
\ifdefempty{\enunciado}{}{%
  \begin{tcolorbox}[enhanced, breakable, colback=acento!4, colframe=acento,
      boxrule=0pt, leftrule=3pt, arc=0pt, title=Enunciado de la tarea,
      coltitle=acento, colbacktitle=acento!4, fonttitle=\bfseries]
    \enunciado
  \end{tcolorbox}
  \vspace{1em}
}

% ---------- Contenido ----------
\section{La situación de aprendizaje}
La situación de aprendizaje es un taller de GeoGebra en el que el alumnado
construye, manipula y observa antes de formalizar, y aplica lo aprendido a
una medición real del centro.

{\small
\renewcommand{\arraystretch}{1.25}
\begin{tabularx}{\linewidth}{>{\bfseries\color{acento}}l X}
\toprule
Título & <<Medir lo que no se puede alcanzar>>: taller de GeoGebra \\
Curso y materia & 1.º de Bachillerato (Ciencias y Tecnología), Matemáticas I \\
Duración & 3 semanas: 12 sesiones de 55 minutos (4 sesiones semanales) \\
Agrupamiento & Equipos heterogéneos de 3--4 alumnos, con trabajo individual en las sesiones de construcción \\
Reto & ¿Cómo podemos medir la altura del edificio del instituto, o la anchura del patio, sin subirnos ni cruzarlo? \\
Producto final & Un libro de GeoGebra por equipo con el modelo de una medición real del centro y una breve exposición oral \\
\bottomrule
\end{tabularx}
}

\section{Marco normativo}
La matriz se construye a partir de la Ley Orgánica 3/2020, de 29 de diciembre
(LOMLOE), que modifica la Ley Orgánica 2/2006, de Educación, y define el
currículo como el conjunto de objetivos, competencias, contenidos, métodos
pedagógicos y criterios de evaluación de cada enseñanza (art.~6). Su
concreción para la Comunidad de Madrid es el Decreto 64/2022, de 20 de julio,
cuyo anexo~II recoge las competencias específicas, los criterios de
evaluación y los contenidos de Matemáticas I. Los contenidos y los criterios
de la matriz se citan de forma literal de este decreto.

Se trabajan siete de las nueve competencias específicas de Matemáticas~I:
modelizar y resolver problemas (CE1), verificar las soluciones, que en el
taller se contrastan con GeoGebra (CE2), conjeturar con ayuda de la
tecnología (CE3), conectar ideas matemáticas (CE5), representarlas con
herramientas digitales (CE7), comunicarlas en la exposición del producto
final (CE8) y trabajar en equipos heterogéneos (CE9).

\section{Cómo leer la matriz}
\begin{itemize}
  \item \textbf{Contenidos:} agrupados por los bloques del decreto según la
    parte del taller en la que se trabajan.
  \item \textbf{Criterio de evaluación:} el criterio de Matemáticas I con el
    que se enlaza cada grupo de contenidos. Un mismo criterio puede aparecer
    en varios grupos, porque se evalúa sobre contenidos distintos.
  \item \textbf{Criterio operativizado:} el criterio redactado de nuevo,
    tomando todo o parte de él y concretándolo en el contenido trabajado, de
    forma que se pueda observar y evaluar.
  \item \textbf{Instrumentos de evaluación:} se marca con una X cada
    instrumento con el que se recoge información sobre ese criterio
    operativizado.
\end{itemize}

\subsection*{Instrumentos de evaluación}
\begin{description}[leftmargin=0pt, itemsep=2pt]
  \item[Examen escrito.] Prueba individual al final de la situación de aprendizaje.
  \item[Cuestionario (GeoGebra Classroom).] Cuestionarios breves de autocorrección al cierre de cada bloque.
  \item[Fichas de construcción en GeoGebra.] Construcciones guiadas de cada sesión de taller, con la pregunta final justificada.
  \item[Libro de GeoGebra.] Producto final del equipo: construcción, cálculos y explicación del método.
  \item[Exposición oral.] Presentación de unos cinco minutos del libro de GeoGebra al resto de la clase.
  \item[Mapa conceptual.] Síntesis en equipo que relaciona vectores, trigonometría y rectas.
  \item[Observación sistemática.] Registro del profesor en las sesiones y seguimiento del tablero Kanban de cada equipo.
  \item[Coevaluación y autoevaluación.] Rúbrica sencilla aplicada a las exposiciones y al trabajo en equipo.
\end{description}

\newgeometry{margin=1.5cm,headheight=14pt}
\begin{landscape}
\section{Matriz de competencias}
\begin{longtblr}[caption={Matriz de competencias de la situación de aprendizaje}, label={tab:matriz}]{
  colspec={Q[l,wd=6.4cm] Q[l,wd=4.8cm] Q[l,wd=6.6cm] Q[c,m,wd=1.5cm] *{8}{Q[c,m,wd=0.55cm]}},
  hlines={0.4pt,gris!50}, vlines={0.4pt,gris!50}, rowhead=2, colsep=3pt,
  row{1-2}={font=\bfseries, fg=white, bg=acento, m}, cells={font=\scriptsize},
}
\SetCell[r=2]{c} Contenidos (Decreto 64/2022) & \SetCell[r=2]{c} Criterio de evaluación & \SetCell[r=2]{c} Criterio operativizado & \SetCell[r=2]{c} Nivel de Bloom & \SetCell[c=8]{c} Instrumentos de evaluación & & & & & & & \\
 & & & & \rotatebox{90}{Examen} & \rotatebox{90}{Cuestionario} & \rotatebox{90}{Fichas GeoGebra} & \rotatebox{90}{Libro GeoGebra} & \rotatebox{90}{Exposición oral} & \rotatebox{90}{Mapa conceptual} & \rotatebox{90}{Observación} & \rotatebox{90}{Co/autoevaluación} \\
\SetCell[r=3]{l,m} {\textbf{A. Números y operaciones: vectores}\\ \textbullet~Adición y producto escalar de vectores: propiedades y representaciones.\\ \textbullet~Conjunto de vectores: estructura, comprensión y propiedades.\\ \textbullet~Módulo de un vector, coordenadas de un vector con respecto a una base, ángulo entre dos vectores y proyección ortogonal.\\ \textbullet~Estrategias para operar con números reales y vectores: cálculo mental o escrito en los casos sencillos y con herramientas tecnológicas en los casos más complicados.} & \textbf{7.1.} Representar ideas matemáticas, estructurando diferentes razonamientos matemáticos y seleccionando las tecnologías más adecuadas. & El alumno ha representado vectores en GeoGebra y ha calculado su suma, su producto por un número y su producto escalar, interpretando geométricamente cada resultado. & \textit{Aplicar} & \textbf{X} & \textbf{X} & \textbf{X} &  &  &  &  &  \\*
 & \textbf{2.1.} Comprobar la validez matemática de las posibles soluciones de un problema utilizando el razonamiento y la argumentación. & El alumno ha comprobado con GeoGebra el módulo, el ángulo entre dos vectores y la proyección ortogonal calculados a mano, justificando el origen de las discrepancias. & \textit{Evaluar} & \textbf{X} &  & \textbf{X} &  &  &  &  &  \\*
 & \textbf{5.2.} Resolver problemas en contextos matemáticos estableciendo y aplicando conexiones entre las diferentes ideas matemáticas. & El alumno ha resuelto problemas de perpendicularidad y de cálculo de ángulos aplicando el producto escalar y relacionando su expresión en coordenadas con su significado geométrico. & \textit{Aplicar} & \textbf{X} & \textbf{X} &  &  &  &  &  &  \\
\SetCell[r=3]{l,m} {\textbf{B y D. Trigonometría}\\ \textbullet~Razones trigonométricas de un ángulo cualquiera medido en grados o en radianes.\\ \textbullet~Demostración de las identidades trigonométricas. Razones trigonométricas del ángulo suma, el ángulo diferencia, el ángulo doble y el ángulo mitad.\\ \textbullet~Cálculo de las razones trigonométricas de un ángulo cualquiera empleando las principales fórmulas trigonométricas.\\ \textbullet~Conjeturas geométricas en el plano: validación por medio de la deducción y la demostración de teoremas.\\ \textbullet~Resolución de ecuaciones trigonométricas sencillas.} & \textbf{3.2.} Emplear herramientas tecnológicas adecuadas en la formulación o investigación de conjeturas o problemas. & El alumno ha examinado en una circunferencia goniométrica dinámica de GeoGebra (con un deslizador para el ángulo) cómo varían el seno, el coseno y la tangente en los cuatro cuadrantes. & \textit{Analizar} &  &  & \textbf{X} &  &  &  & \textbf{X} &  \\*
 & \textbf{3.1.} Adquirir nuevo conocimiento matemático a partir de la formulación de conjeturas y problemas de forma guiada. & El alumno ha formulado, de forma guiada, conjeturas sobre identidades trigonométricas (fundamental, ángulo doble, ángulo suma) a partir de la construcción y las ha demostrado por escrito. & \textit{Crear} & \textbf{X} &  & \textbf{X} &  &  &  &  &  \\*
 & \textbf{1.2.} Obtener todas las posibles soluciones matemáticas de problemas de la vida cotidiana y de la ciencia y la tecnología, describiendo el procedimiento utilizado. & El alumno ha resuelto ecuaciones trigonométricas sencillas obteniendo todas sus soluciones, tanto en [0°, 360°) como en forma general, y ha descrito el procedimiento seguido. & \textit{Aplicar} & \textbf{X} & \textbf{X} &  &  &  &  &  &  \\
\SetCell[r=3]{l,m} {\textbf{B. Medida y geometría: resolución de triángulos}\\ \textbullet~Cálculo de longitudes y medidas angulares: uso de la trigonometría.\\ \textbullet~Aplicación de las razones trigonométricas, el teorema de los senos y el teorema del coseno en la resolución de triángulos y de problemas geométricos de contexto real.\\ \textbullet~Demostración del teorema del seno y del coseno.} & \textbf{1.1.} Manejar algunas estrategias y herramientas, incluidas las digitales, en la modelización y resolución de problemas de la vida cotidiana y de la ciencia y la tecnología, evaluando su eficiencia en cada caso. & El alumno ha justificado la elección entre el teorema del seno y el del coseno según los datos de cada triángulo, valorando la eficiencia de cada estrategia. & \textit{Evaluar} & \textbf{X} &  & \textbf{X} &  &  &  &  &  \\*
 & \textbf{1.2.} Obtener todas las posibles soluciones matemáticas de problemas de la vida cotidiana y de la ciencia y la tecnología, describiendo el procedimiento utilizado. & El alumno ha discriminado si un triángulo dado por dos lados y el ángulo opuesto a uno de ellos tiene dos, una o ninguna solución y ha calculado todas las soluciones, describiendo el procedimiento. & \textit{Analizar} & \textbf{X} &  & \textbf{X} &  &  &  &  &  \\*
 & \textbf{2.1.} Comprobar la validez matemática de las posibles soluciones de un problema utilizando el razonamiento y la argumentación. & El alumno ha verificado en GeoGebra, arrastrando los vértices, la validez de los triángulos obtenidos analíticamente y ha argumentado por qué un caso no tiene solución. & \textit{Evaluar} &  & \textbf{X} & \textbf{X} &  &  &  &  &  \\
\SetCell[r=3]{l,m} {\textbf{C. Geometría en el plano: rectas}\\ \textbullet~Ecuaciones de la recta en el espacio bidimensional.\\ \textbullet~Estudio de la posición relativa de puntos y rectas en el plano.\\ \textbullet~Planteamiento y resolución de problemas de geometría afín relacionados con la incidencia, el paralelismo y la ortogonalidad de rectas en el plano.\\ \textbullet~Planteamiento y resolución de problemas de geometría métrica relacionados con la medida de ángulos entre rectas y la medida de distancias entre puntos y rectas.\\ \textbullet~Expresiones algebraicas de objetos geométricos: selección de la más adecuada en función de la situación a resolver.} & \textbf{5.1.} Manifestar una visión matemática integrada, investigando y conectando las diferentes ideas matemáticas. & El alumno ha relacionado el vector director y el vector normal con las distintas ecuaciones de la recta, organizando en un mapa conceptual las conexiones entre vectores, trigonometría y geometría analítica. & \textit{Analizar} & \textbf{X} &  &  &  &  & \textbf{X} &  &  \\*
 & \textbf{5.2.} Resolver problemas en contextos matemáticos estableciendo y aplicando conexiones entre las diferentes ideas matemáticas. & El alumno ha determinado la posición relativa de dos rectas y ha calculado ángulos entre rectas y distancias de un punto a una recta aplicando el producto escalar. & \textit{Aplicar} & \textbf{X} & \textbf{X} &  &  &  &  &  &  \\*
 & \textbf{7.2.} Seleccionar y utilizar diversas formas de representación, valorando su utilidad para compartir información. & El alumno ha seleccionado la ecuación de la recta (vectorial, paramétrica, continua, general o explícita) más útil para cada problema, justificando su elección. & \textit{Evaluar} & \textbf{X} &  & \textbf{X} &  &  &  &  &  \\
\SetCell[r=4]{l,m} {\textbf{C y D. Modelización: producto final «Medir lo que no se puede alcanzar»}\\ \textbullet~Representación de objetos geométricos en el plano mediante herramientas digitales.\\ \textbullet~Modelos matemáticos (geométricos, algebraicos, grafos...) en la resolución de problemas en el plano. Conexiones con otras disciplinas y áreas de interés.\\ \textbullet~Modelización de la posición y el movimiento de un objeto en el plano mediante vectores.\\ \textbullet~Formulación, resolución y análisis de problemas de la vida cotidiana y de la ciencia y la tecnología utilizando herramientas o programas adecuados.} & \textbf{1.1.} Manejar algunas estrategias y herramientas, incluidas las digitales, en la modelización y resolución de problemas de la vida cotidiana y de la ciencia y la tecnología, evaluando su eficiencia en cada caso. & El alumno ha diseñado un método de medida y ha construido un modelo a escala en GeoGebra para calcular una altura o distancia inaccesible del centro, valorando su eficiencia frente a otros métodos. & \textit{Crear} &  &  &  & \textbf{X} & \textbf{X} &  &  &  \\*
 & \textbf{2.2.} Seleccionar la solución más adecuada de un problema en función del contexto, usando el razonamiento y la argumentación. & El alumno ha valorado el efecto de un error de un grado en el resultado final y ha seleccionado, entre las medidas repetidas, el resultado más razonable. & \textit{Evaluar} &  &  &  & \textbf{X} & \textbf{X} &  &  &  \\*
 & \textbf{7.2.} Seleccionar y utilizar diversas formas de representación, valorando su utilidad para compartir información. & El alumno ha elaborado un libro de GeoGebra que integra la construcción, los cálculos y la explicación del método, eligiendo las representaciones más claras para compartirlo con la clase. & \textit{Crear} &  &  &  & \textbf{X} &  &  &  & \textbf{X} \\*
 & \textbf{8.1.} Mostrar organización al comunicar las ideas matemáticas empleando el soporte, la terminología y el rigor apropiados. & El alumno ha explicado oralmente el método y los resultados con orden, utilizando la terminología y la notación matemáticas con rigor. & \textit{Comprender} &  &  &  &  & \textbf{X} &  &  & \textbf{X} \\
\SetCell[r=3]{l,m} {\textbf{F. Actitudes y aprendizaje}\\ \textbullet~Tratamiento del error como elemento movilizador de conocimientos previos adquiridos y generador de oportunidades de aprendizaje en el aula de matemáticas.\\ \textbullet~Reconocimiento y aceptación de diversos planteamientos en la resolución de problemas, transformando los enfoques de los demás en nuevas y mejoradas estrategias propias.\\ \textbullet~Técnicas y estrategias de trabajo en equipo para la resolución de problemas y tareas matemáticas, en grupos heterogéneos.\\ \textbullet~Destrezas para desarrollar una comunicación efectiva: la escucha activa, la formulación de preguntas o solicitud y prestación de ayuda cuando sea necesario.} & \textbf{9.1.} Afrontar las situaciones de incertidumbre, aceptando y aprendiendo del error como parte del proceso de aprendizaje de las matemáticas. & El alumno ha analizado los errores de sus conjeturas y de sus medidas, y los ha utilizado como punto de partida para revisar su razonamiento en la puesta en común. & \textit{Analizar} &  &  &  &  &  &  & \textbf{X} & \textbf{X} \\*
 & \textbf{9.2.} Mostrar una actitud positiva y perseverante, aceptando y aprendiendo de la crítica razonada al hacer frente a las diferentes actividades de las matemáticas. & El alumno ha valorado la crítica razonada recibida en la coevaluación y la ha incorporado para mejorar el libro de GeoGebra, perseverando ante las dificultades. & \textit{Evaluar} &  &  &  &  &  &  & \textbf{X} & \textbf{X} \\*
 & \textbf{9.3.} Participar en tareas matemáticas de forma activa en equipos heterogéneos, escuchando su razonamiento. & El alumno ha colaborado activamente en el equipo heterogéneo, ha asumido tareas del tablero Kanban y ha tenido en cuenta los razonamientos de sus compañeros. & \textit{Aplicar} &  &  &  &  &  &  & \textbf{X} & \textbf{X} \\
\end{longtblr}

\end{landscape}
\restoregeometry

\end{document}
